\documentclass[a4paper,11pt]{article}
\usepackage{iftex}
\iftutex
  % Si compilation avec LuaLaTeX ou XeLaTeX
  \usepackage{fontspec}
\else
  % Si compilation avec pdfLaTeX
  \usepackage[utf8]{inputenc}
  \usepackage[T1]{fontenc}
\fi
\usepackage[french]{babel}
\usepackage{geometry}
\geometry{hmargin=2.5cm, vmargin=2.5cm}
\usepackage{hyperref}
\usepackage{xcolor}
\usepackage{listings}

\hypersetup{
    colorlinks=true,
    linkcolor=blue,
    filecolor=magenta,      
    urlcolor=cyan,
}

\lstset{
  language=[LaTeX]TeX,
  basicstyle=\ttfamily\small,
  keywordstyle=\color{blue},
  commentstyle=\color{gray},
  frame=single,
  backgroundcolor=\color{black!5},
  breaklines=true,
  postbreak=\mbox{\textcolor{gray}{$\hookrightarrow$}\space}
}

\title{\textbf{Le Paquet \texttt{archery-trainer-sets}}}
\author{deogracia}
\newcommand{\pkgversion}{1.1.0-dev}
\date{\today \quad -- \quad Version \pkgversion}

\begin{document}
\maketitle

\begin{abstract}
Le paquet \texttt{archery-trainer-sets} regroupe un ensemble de classes et de paquets \LaTeX{} destinés aux entraîneurs de tir à l'arc. Il permet de standardiser et de faciliter la création de documents utiles à la préparation et au suivi des séances d'entraînement.
\end{abstract}

\tableofcontents
\vspace{1cm}

\section{Introduction}
Ce paquet a été conçu pour répondre aux besoins spécifiques des entraîneurs de tir à l'arc (initiation, progression, compétition). Il inclut pour le moment la classe \texttt{ficheseance}, qui permet de générer des fiches de séance claires, structurées et prêtes à être imprimées (format Paysage, 7 colonnes).

Le paquet est distribué sous licence \textbf{LPPL 1.3c}.

\section{La classe \texttt{ficheseance}}

\subsection{Chargement de la classe}
La classe \texttt{ficheseance} est basée sur la classe standard \texttt{article}. Elle configure automatiquement le document en format A4, orientation paysage, avec une police de 10pt. 
Pour l'utiliser, il suffit de déclarer :
\begin{lstlisting}
\documentclass{ficheseance}
\end{lstlisting}

La classe gère nativement le français et l'anglais via le paquet \texttt{translations}. Il suffit de charger \texttt{babel} ou \texttt{polyglossia} avec la langue de votre choix pour que les en-têtes se traduisent automatiquement.

\subsection{Définition des métadonnées}
Avant le \verb|\begin{document}| ou juste après, vous devez définir les informations de votre séance à l'aide des commandes suivantes :
\begin{itemize}
    \item \verb|\numfiche{...}| : Numéro de la fiche ou de la séance.
    \item \verb|\datefiche{...}| : Date de la séance.
    \item \verb|\archer{...}| : Nom de l'archer ou du groupe.
    \item \verb|\theme{...}| : Thème général de la séance.
    \item \verb|\objectifseance{...}| : Objectif principal de la séance.
\end{itemize}

\subsection{Génération de l'en-tête}
Une fois les informations définies, utilisez la commande \verb|\makeficheheader| pour afficher le cartouche récapitulatif en haut de la première page.

\subsection{Le tableau de séance}
Le cœur de la fiche est l'environnement \texttt{tableauSeance}. Il déploie un tableau multipage (grâce à \texttt{xltabular}) composé de 7 colonnes spécifiques à l'entraînement sportif :
\begin{enumerate}
    \item \textbf{Objectif}
    \item \textbf{Temps}
    \item \textbf{Situations}
    \item \textbf{Outils}
    \item \textbf{Consignes}
    \item \textbf{Critères de réussite}
    \item \textbf{Évaluations}
\end{enumerate}

Pour insérer une ligne (une situation ou un exercice) dans le tableau, utilisez la commande \verb|\ligneSeance| qui prend exactement 7 arguments :
\begin{lstlisting}
\begin{tableauSeance}
    \ligneSeance
        {Échauffement} % Objectif
        {15 min} % Temps
        {Réveils articulaires et musculaires} % Situation
        {Élastiques} % Outils
        {Mouvements lents et contrôlés} % Consignes
        {Pas de douleur} % Critères
        {OK} % Évaluation
\end{tableauSeance}
\end{lstlisting}

\section{Support et Contribution}
Le code source et les exemples complets sont disponibles sur le dépôt GitHub :\\ \url{https://github.com/deogracia/archery-trainer-sets}.\\
Les contributions sous forme de \textit{Pull Requests} sont les bienvenues.

\end{document}